\section{Заключение}
\label{sec:conclusion}

Исследование показало, что для управления импульсным возбуждением КТ
вблизи золотой МНЧ необходимо совместно выбирать место расположения
КТ, поляризацию поля и рабочую частоту. Для рассмотренного сфероида
с полуосями \SI{10}{\nano\meter} и \SI{15}{\nano\meter} и импульса
длительностью \SI{20}{\femto\second} наиболее эффективна торцевая
конфигурация с полем вдоль оси симметрии. При поверхностном зазоре
\SI{1}{\nano\meter} флюенс достижения населённости $0.5$ составляет
$1.62\cdot10^{-6}$~Дж/см$^2$ против
$6.167\cdot10^{-5}$~Дж/см$^2$ для изолированной КТ.
Полученный в модели выигрыш примерно в $38$ раз делает эту
конфигурацию предпочтительным кандидатом для экспериментальной
проверки. Боковое радиальное возбуждение также даёт выигрыш,
тогда как смена поляризации может приводить к сильному подавлению
отклика. Эффект определяется всем связанным откликом, а не только
расстоянием или ориентацией поля относительно локальной поверхности.

Пространственное уточнение поля и уточнение материальной дисперсии
оказались существенными для количественных выводов. В ведущем канале
при зазоре \SI{1}{\nano\meter} дипольное приближение завышает порог
в $2.4$ раза; согласие на уровне $10\,\%$ на исследованной сетке
в пределах принятого квазистатического критерия не достигается.
Для нескольких других каналов такое согласие устанавливается
при зазорах $7$--$\SI{10}{\nano\meter}$. Одноосцилляторная
аппроксимация выбранной дисперсии золота даёт в торцевом
сценарии при $g=\SI{0.5}{\nano\meter}$ порог в $7.7$ раза выше
многоосцилляторной. Следовательно,
ни точечно-дипольную замену МНЧ, ни описание её отклика одним
осциллятором нельзя принимать без проверки в рабочей геометрии
и спектральной полосе.

Практический выбор зазора зависит от измеряемой величины.
Слабополевое возбуждение на фиксированной частоте имеет максимум
вблизи \SI{1.5}{\nano\meter}, возникающий из конкуренции усиления
поля и сдвига резонанса. Импульсный порог уменьшается до нижней
границы расчётной сетки, но этот результат не обосновывает
требование максимально сблизить реальные частицы. Диапазон
$1$--$\SI{1.5}{\nano\meter}$ предложен как отправная точка
для сравнения с более толстыми разделителями. Его пригодность
определяется сохранением возбуждения и допустимыми потерями,
а не только минимальным расчётным флюенсом.

Для получения образца наиболее обоснован маршрут, сочетающий
химически селективное присоединение КТ к концам стержней
и последующую ориентацию комплексов в прозрачной матрице.
Спектроскопическая проверка должна объединять поляризационную
экстинкцию, КТ-селективные спектры, зависимости от флюенса
и измерение релаксации. Совместная обработка этих данных с учётом
распределения параметров позволяет проверять гипотезу
о преобладании торцевой или боковой конфигурации. Один спектр
люминесценции не устанавливает численные доли таких групп,
поскольку вклад каждой пары зависит также от её квантового
выхода и эффективности регистрации.

Сформулированные рекомендации относятся к принятой полуклассической
модели с фиксированными скоростями релаксации. Конечный размер КТ,
геометрически зависимый спонтанный и безызлучательный распад,
подложка и взаимодействие соседних пар могут изменить абсолютные
пороги и их соотношения. Поэтому основной итог работы --
количественно обоснованный внутри этой модели выбор геометрии
для дальнейшего изготовления и набор измерений, позволяющих
проверить её преимущество. Переход от такого выбора к подтверждённым
характеристикам композита требует согласования расчёта
с измеренными спектрами и временами релаксации конкретного образца.
